\documentclass[12pt]{article}
\usepackage{amsmath, amssymb}
\usepackage{geometry}
\geometry{margin=1in}

\title{CSE400 -- Fundamentals of Probability in Computing\\
Lecture L11\_12: Transformation of Random Variables}
\date{}
\begin{document}
\maketitle

\section*{1. Transformation of Random Variables}

\subsection*{1.1 Definition}

Let $X$ be a random variable with known distribution.

Let
\[
Y = g(X)
\]

Objective: Find the distribution (PDF/PMF) of $Y$.

\subsection*{1.2 Transformation -- Discrete Case}

If $X$ is discrete with PMF $p_X(x)$, and
\[
Y = g(X),
\]
then
\[
p_Y(y) = P(Y = y)
\]

Since $Y = g(X)$,
\[
p_Y(y) = P(g(X) = y)
\]

Thus,
\[
p_Y(y) = \sum_{x: g(x)=y} p_X(x)
\]

\[
\boxed{
p_Y(y) = \sum_{x: g(x)=y} p_X(x)
}
\]

\subsection*{1.3 Transformation -- Continuous Case}

Let $X$ be continuous with PDF $f_X(x)$.

Let
\[
Y = g(X)
\]

\subsubsection*{Case 1: Monotonic Transformation}

Assume:
\begin{itemize}
    \item $g(\cdot)$ is one-to-one and differentiable.
    \item Inverse exists: $x = g^{-1}(y)$.
\end{itemize}

\textbf{Step 1: CDF Method}
\[
F_Y(y) = P(Y \le y)
\]

If $g$ is increasing:
\[
F_Y(y) = P(g(X) \le y)
\]
\[
= P(X \le g^{-1}(y))
\]
\[
= F_X(g^{-1}(y))
\]

\textbf{Step 2: Differentiate to get PDF}
\[
f_Y(y) = \frac{d}{dy} F_Y(y)
\]
\[
= \frac{d}{dy} F_X(g^{-1}(y))
\]

Using chain rule:
\[
f_Y(y) = f_X(g^{-1}(y)) \cdot \frac{d}{dy} g^{-1}(y)
\]

If decreasing, absolute value is required.

\[
\boxed{
f_Y(y) =
f_X(g^{-1}(y))
\left|
\frac{d}{dy} g^{-1}(y)
\right|
}
\]

\subsection*{1.4 Multiple Roots Case}

If $g(x)$ is not one-to-one and
\[
y = g(x)
\]
has multiple solutions $x_1, x_2, \dots, x_n$, then

\[
\boxed{
f_Y(y)
=
\sum_i
f_X(x_i)
\left|
\frac{dx_i}{dy}
\right|
}
\]

\section*{2. Function of Two Random Variables}

Let $X, Y$ have joint PDF:
\[
f_{X,Y}(x,y)
\]

Define:
\[
U = g_1(X,Y)
\]
\[
V = g_2(X,Y)
\]

We want:
\[
f_{U,V}(u,v)
\]

\subsection*{2.1 One-to-One Transformation}

Assume:
\[
x = x(u,v)
\]
\[
y = y(u,v)
\]

\subsection*{2.2 Jacobian Determinant}

\[
J =
\begin{vmatrix}
\frac{\partial x}{\partial u} & \frac{\partial x}{\partial v} \\
\frac{\partial y}{\partial u} & \frac{\partial y}{\partial v}
\end{vmatrix}
\]

\[
J =
\frac{\partial x}{\partial u}\frac{\partial y}{\partial v}
-
\frac{\partial x}{\partial v}\frac{\partial y}{\partial u}
\]

\subsection*{2.3 Joint PDF Formula}

\[
\boxed{
f_{U,V}(u,v)
=
f_{X,Y}(x(u,v), y(u,v))
\cdot
|J|
}
\]

\section*{3. Illustrative Example: $Z = X + Y$}

Given two continuous random variables $X$ and $Y$.

Define:
\[
Z = X + Y
\]

\subsection*{3.1 Define Transformation}

\[
U = X + Y
\]
\[
V = Y
\]

Then:
\[
X = U - V
\]
\[
Y = V
\]

\subsection*{3.2 Jacobian}

\[
J =
\begin{vmatrix}
\frac{\partial X}{\partial U} & \frac{\partial X}{\partial V} \\
\frac{\partial Y}{\partial U} & \frac{\partial Y}{\partial V}
\end{vmatrix}
\]

\[
\frac{\partial X}{\partial U} = 1
\quad
\frac{\partial X}{\partial V} = -1
\]
\[
\frac{\partial Y}{\partial U} = 0
\quad
\frac{\partial Y}{\partial V} = 1
\]

\[
J =
\begin{vmatrix}
1 & -1 \\
0 & 1
\end{vmatrix}
\]

\[
J = (1)(1) - (-1)(0) = 1
\]

\[
|J| = 1
\]

\subsection*{3.3 Joint PDF of $(U,V)$}

\[
f_{U,V}(u,v)
=
f_{X,Y}(u-v, v)
\]

\subsection*{3.4 Marginal PDF of $Z = U$}

\[
f_Z(u)
=
\int_{-\infty}^{\infty}
f_{U,V}(u,v) \, dv
\]

\[
=
\int_{-\infty}^{\infty}
f_{X,Y}(u-v, v) \, dv
\]

\[
\boxed{
f_Z(z)
=
\int_{-\infty}^{\infty}
f_{X,Y}(z-y, y) \, dy
}
\]

\subsection*{3.5 If $X$ and $Y$ are Independent}

\[
f_{X,Y}(x,y)
=
f_X(x) f_Y(y)
\]

\[
f_Z(z)
=
\int_{-\infty}^{\infty}
f_X(z-y) f_Y(y) \, dy
\]

\[
\boxed{
f_Z(z)
=
\int_{-\infty}^{\infty}
f_X(z-y) f_Y(y) \, dy
}
\]

This is the convolution of $f_X$ and $f_Y$.

\section*{Final Summary (Exam Ready)}

\[
p_Y(y)
=
\sum_{x: g(x)=y} p_X(x)
\]

\[
f_Y(y)
=
f_X(g^{-1}(y))
\left|
\frac{d}{dy} g^{-1}(y)
\right|
\]

\[
f_Y(y)
=
\sum_i
f_X(x_i)
\left|
\frac{dx_i}{dy}
\right|
\]

\[
f_{U,V}(u,v)
=
f_{X,Y}(x(u,v), y(u,v))
\cdot |J|
\]

\[
f_Z(z)
=
\int_{-\infty}^{\infty}
f_X(z-y) f_Y(y) dy
\]

\end{document}
